\minitopic{数字证据不仅有来源，还有时间}网页、数据库和生成系统持续更新；今天打开的内容可能已经不是使用者昨天据以行动的版本。只保存链接不能恢复认识过程，因为同一地址下的正文、模型和检索排序都可能改变。核验记录应同时保存访问时间、内容版本、关键摘录和必要的存档；对会持续变化的规范、产品能力和公共通知，还要标明结论的有效截止时间。

时间也改变证据权重。较新的材料不必然更可靠：突发事件早期消息更新快却未经核实，较旧的正式报告可能证据更完整；旧材料也不必然失效，基础定义和历史记录可以长期稳定。读者应问变化发生在世界、数据、解释还是发布状态。把“最新”当成“最好”，会奖励速度；忽略更新时间，则会把已经撤回或修订的主张继续传播。

\minitopic{更正必须沿传播链返回}数字内容可以在数分钟内被复制到截图、聚合页和生成回答中，原页面的更正不会自动传播。一个负责任系统应保留主张的来源标识，并在重要来源撤回、改版或发生冲突时，能够找到受影响的下游内容。否则系统只擅长把信息向前分发，却没有把反证向后送达的能力。

这形成一种认识债务：系统越容易复用旧输出，未来核查和更正的负担越大。降低债务的方法包括为高风险主张设置失效日期、区分原始材料与缓存摘要、记录哪些结论依赖哪些版本，并让使用者订阅重要更正。真正可纠错的界面不仅在回答时显示来源，还在来源变化时承认过去答案需要重审。

\begin{argumentbox}
对时间敏感的数字主张，最小记录不是“链接+截图”，而是“精确命题+获取时间+来源版本+支持段落+当时系统版本+复查日期”。这组信息让后来者能够区分：原判断当时是否合理，以及它现在是否仍然成立。
\end{argumentbox}

\index{版本责任}
\index{认识债务}
